\begin{afterword}
\summaryhead

At the 2007 Singularity Summit a speaker called for democratic, multinational development of artificial intelligence. The author asked which would be more democratic: a consortium of democracies full of ordinary politicking, or rebel nerds whose AI polls the world and does what the majority says. The speaker called the first an editorial in \textsc{Reason} and the second a Hollywood plot, offered the Human Genome Project as a model, and, asked how interest groups would resolve their conflicts in such a structure, said ``I don't know.''

The author argues that the substance of democracy is a mechanism for resolving policy conflicts. To call for democracy without one is to say a word the audience is meant to cheer, like the ``Applause'' sign in a television studio. Many such applause lights can be caught by reversing them: if the reversed sentence sounds abnormal, the original conveys no new information. Such sentences have legitimate uses, but if no specifics follow, the sentence is probably an applause light. The post ends with a sample talk made of nothing else.

\responsehead

The tool the post offers is useful. Reversing a sentence to see whether anyone would say the reverse is a quick check on empty speech, and the post limits it sensibly: an uninformative sentence can introduce a topic or stress a proposal, and becomes suspect only if no specifics follow. The underlying observation is Orwell's, that ``democracy'' is used as praise with no agreed definition. The author said in the comments that Orwell was in mind while writing; the post does not mention him.

The example is weaker than the tool. It is one exchange, told from memory, with the speaker unnamed. After the speaker, Jamais Cascio, objected in the comments, the author, pointing readers to the recording, wrote that the event ``did not occur the exact way I remembered it'': the Genome Project example had been volunteered, not extracted by a follow-up question. The text as published today still presents it as extracted. Cascio's account also supplies what the post leaves out of two replies: that both of the author's scenarios were caricatures with benefits and drawbacks, and that ``I don't know'' was followed by ``this is not a solved problem, but it's an important problem.''

Once those parts are put back, little supports the diagnosis. The post concludes, with ``I think,'' that the word was said so the audience would cheer. What the exchange shows is that the speaker had not worked out how an international AI project would settle conflicts between interest groups, and said so. That is an open question, not a word used for applause, and by the post's own rule a specific did follow. The author granted as much in the thread (``I do not accuse you too much of this'').

In short: a useful test for empty speech, which goes back to Orwell, who is not credited; the test is illustrated with a remembered exchange that the author later found had gone differently, and whose speaker met the post's own condition by naming a specific model.
\end{afterword}
